\section{Theoretical Foundations and Proofs}
\label{app:theory}

This appendix develops the compact theory used in the main text. We begin with the pointwise decomposition, then lift it to a trajectory, a population, and a decision problem. The proofs require only integrability of the signed response and do not assume differentiability, smoothness, or a particular model class.

\subsection{Positive and negative parts}
\label{app:jordan}

For a real number $u$, define $u^{+}=\max\{u,0\}$ and $u^{-}=\max\{-u,0\}$. Then
\begin{equation}
    u=u^{+}-u^{-},
    \qquad
    |u|=u^{+}+u^{-},
    \qquad
    u^{+}u^{-}=0.
    \label{eq:positive-negative}
\end{equation}
For an integrable function $z:\mathcal T\rightarrow\mathbb R$, the signed measure $\nu(A)=\int_A z(t)\,\mathrm d\mu(t)$ admits the Jordan decomposition $\nu=\nu^{+}-\nu^{-}$, where $\mathrm d\nu^{+}=z^{+}\mathrm d\mu$ and $\mathrm d\nu^{-}=z^{-}\mathrm d\mu$. DECAF applies this decomposition to the endpoint-oriented response on the active branch and reserves a separate branch for endpoint-null pairs.

\subsection{Canonicality}

\begin{proof}[Proof of Theorem~\ref{thm:canonical}]
Fix $a$, $s$, and $r$. If $a=0$, endpoint gating requires $e=c=0$. Conservation then forces $f=|r|$. Hence the triple is unique.

Now suppose $a=1$. Endpoint gating gives $f=0$. Let $z=sr$. If $z\geq0$, directional support requires $c=0$ and conservation gives $e=|r|=z=z^{+}$. If $z\leq0$, directional support requires $e=0$ and conservation gives $c=|r|=-z=z^{-}$. Thus
\begin{equation}
    (e,c,f)=\bigl(a(sr)^{+},a(sr)^{-},(1-a)|r|\bigr).
\end{equation}
The construction is nonnegative, conserves $|r|$, obeys the gate, and has the stated support. Therefore it is the unique triple satisfying the axioms.
\end{proof}

The theorem also yields a projection interpretation. On an active pair, $e$ is the magnitude of the Euclidean projection of $r$ onto the endpoint-aligned ray $\{s\lambda:\lambda\geq0\}$, while $c$ is the magnitude of the projection onto the opposite ray. The null branch is not a projection. It records that the endpoint does not provide a stable orientation at the chosen threshold.

\subsection{Magnitude non-identifiability and strict refinement}

\begin{proof}[Proof of Theorem~\ref{thm:nonident}]
Fix $m>0$. Consider three responses. First, let the endpoint be active with $s=1$ and let $r=m$. Then $(e,c,f)=(m,0,0)$. Second, keep the same endpoint and let $r=-m$. Then $(e,c,f)=(0,m,0)$. Third, let the endpoint be null and let $r=m$. Then $(e,c,f)=(0,0,m)$. In every case, the observed magnitude is $|r|=m$. A statistic that observes only $|r|$ therefore cannot distinguish the three mechanisms.
\end{proof}

\begin{corollary}[Strict response refinement]
\label{cor:strict-refinement}
The DECAF profile is at least as informative as ordinary magnitude for every decision problem, because $\Abs=E+C+F$ is a deterministic function of the profile. It is strictly more informative for any mechanism-identification problem that assigns positive probability to two distinct profiles with the same magnitude.
\end{corollary}

\begin{proof}
Any rule that uses $\Abs$ can be composed with the map $(E,C,F)\mapsto E+C+F$. Theorem~\ref{thm:nonident} gives distinct profiles that share one magnitude. A decision problem that rewards correct mechanism identification separates those profiles, so no magnitude-only rule can match the best rule that observes the full profile. This is the standard logic of strict Blackwell refinement~\citep{blackwell1953equivalent}.
\end{proof}

A useful finite case makes the gap concrete. Suppose evidence, contradiction, and fragility are equally likely and all produce magnitude $m$. Every magnitude-only classifier receives the same observation, so its best accuracy is $1/3$. Under unequal priors, the best magnitude-only accuracy is the largest prior. The DECAF profile identifies the mechanism exactly in this idealized construction.

\subsection{Preservation, attenuation, and inversion}

\begin{proof}[Proof of Proposition~\ref{prop:calibration}]
Let the endpoint orientation be $s$. Under preservation, the stage response is $sm$, so $E=m$, $C=0$, and $\Abs=m$.

Under attenuation, the response is $sm$ with probability $1-\eta$ and zero with probability $\eta$. Taking expectations gives $E=(1-\eta)m$, $C=0$, and $\Abs=(1-\eta)m$.

Under inversion, the response is $sm$ with probability $1-\eta$ and $-sm$ with probability $\eta$. The first branch contributes $m$ to evidence, while the second contributes $m$ to contradiction. Therefore $E=(1-\eta)m$, $C=\eta m$, and $\Abs=m$. The ratio is $C/(E+C)=\eta$.
\end{proof}

The equal-magnitude assumption isolates the inversion probability. If the aligned and opposed branches have magnitudes $m_{+}$ and $m_{-}$, then
\begin{equation}
    E=(1-\eta)m_{+},
    \qquad
    C=\eta m_{-},
    \qquad
    \frac{C}{E+C}=\frac{\eta m_{-}}{(1-\eta)m_{+}+\eta m_{-}}.
    \label{eq:unequal-inversion}
\end{equation}
Thus the opposed fraction measures probability mass only when the two branches have comparable effect size. In general, $C$ measures opposed causal mass.


\subsection{Population factorization}
\label{app:population-factorization}

Let $\pi=\Pr(|d|\geq\varepsilon)$. Define branch-conditional summaries whenever the corresponding branch has positive probability. Since evidence and contradiction vanish on the null branch, while fragility vanishes on the active branch,
\begin{equation}
    E=\pi E^{\mathrm{active}},
    \qquad
    C=\pi C^{\mathrm{active}},
    \qquad
    F=(1-\pi)F^{\mathrm{null}}.
    \label{eq:population-factorization}
\end{equation}
Unconditional quantities combine branch prevalence and conditional intensity. Conditional quantities answer a different question and become unstable when their branch is rare. The factorization is therefore useful when an external outcome is defined only within the active or null branch, but the primary experiments report unconditional population summaries.

\subsection{Invariances}
\label{app:invariances}

\paragraph{Endpoint swap.}
Swapping $\mathbf x^{+}$ and $\mathbf x^{-}$ multiplies both $d$ and $r(t)$ by $-1$. The oriented response $s r(t)$ is unchanged. Hence $M,E,C,F,$ and $\Abs$ are invariant.

\paragraph{Score translation.}
Replacing $q$ by $q+b$ leaves every difference unchanged.

\paragraph{Positive affine score reparameterization.}
Let $q'=\lambda q+b$ with $\lambda>0$. Then
$d'=\lambda d$, $r'(t)=\lambda r(t)$, and $s'=s$.
If the endpoint threshold is transformed consistently as
$\varepsilon'=\lambda\varepsilon$, the gate is unchanged and
\[
(M',E',C',F',\Abs')
=
\lambda(M,E,C,F,\Abs).
\]
Hence all normalized component proportions are invariant.

With a numerically fixed threshold $\varepsilon$, however, the gate obeys
\[
a'_{\varepsilon}
=
\mathbf{1}\{|d|\ge\varepsilon/\lambda\}.
\]
Consequently,
\[
E'_{\varepsilon}=\lambda E_{\varepsilon/\lambda},\qquad
C'_{\varepsilon}=\lambda C_{\varepsilon/\lambda},\qquad
F'_{\varepsilon}=\lambda F_{\varepsilon/\lambda}.
\]
Thus positive scaling preserves component proportions only when gate
membership is preserved. Otherwise, the discrepancy is controlled by
the threshold-mass bound in Appendix~B.8.

\paragraph{Negative score scaling.}
A negative scaling reverses the semantic meaning of the target score. The endpoint orientation still makes the component magnitudes invariant after multiplying by $|\lambda|$, but the analyst has changed the behavior being explained. We therefore define the score direction before analysis.

\subsection{Protocol reparameterization}
\label{app:reparameterization}

Let $h:\widetilde{\mathcal T}\rightarrow\mathcal T$ be an increasing bijection and define the reparameterized path $\widetilde r(u)=r(h(u))$. Pointwise components are reindexed:
\begin{equation}
    \widetilde e(u)=e(h(u)),
    \quad
    \widetilde c(u)=c(h(u)),
    \quad
    \widetilde f(u)=f(h(u)).
\end{equation}
If the integration measure is pushed forward consistently, $\widetilde\mu=h^{-1}_{\#}\mu$, then the integrated profile is invariant. If both paths are instead integrated with a uniform coordinate measure, their AUCs may differ. Dynamic DECAF summaries are therefore protocol-relative, and every experiment reports the path and stage measure.

\subsection{Threshold stability}
\label{app:threshold-theory}

Let $0\leq\varepsilon_1<\varepsilon_2$ and suppose the per-example path summary is bounded by $B$. The two gates differ only on pairs satisfying $\varepsilon_1\leq|d|<\varepsilon_2$. Consequently, for each component $G\in\{E,C,F\}$,
\begin{equation}
    |G_{\varepsilon_2}-G_{\varepsilon_1}|
    \leq
    B\,\Pr\!\left(\varepsilon_1\leq|d|<\varepsilon_2\right).
    \label{eq:threshold-bound}
\end{equation}
The threshold is stable when little endpoint mass lies near the boundary. We report sensitivity over multiple thresholds rather than treating one numerical cutoff as universal.
